
Prior work has studied verification strategies for LLM code generation in isolation.

\subsubsection{2.1 Grammar-Constrained Decoding}

Grammar-constrained decoding enforces syntactic validity during generation by masking logits that would produce invalid tokens. Mündler et al. (2025) introduced type-constrained decoding using prefix automata, reducing compilation errors by over 50\% on HumanEval and MBPP. SynCode extends this to general-purpose programming languages with soundness and completeness guarantees. CRANE \citep{banerjee2025crane} adds reasoning augmentation to constrained decoding, achieving 10\% accuracy improvement on GSM-symbolic and FOLIO.

These approaches operate at the token level during generation, preventing syntax errors before they occur. However, they do not address semantic issues or specification violations.

\subsubsection{2.2 Static Analysis Feedback Loops}

Post-generation static analysis identifies semantic patterns that grammar constraints cannot catch. Blyth et al. (2025) demonstrated that static analysis-driven prompting reduces security issues from over 40\% to 13\% and reliability issues from over 50\% to 11\% on PythonSecurityEval. PropertyGPT \citep{liu2024propertygpt} combines retrieval-augmented generation with static analysis feedback for smart contract property generation, achieving 80\% recall and detecting 26 CVEs.

These approaches rely on the LLM using feedback to repair code. The present work reveals that this assumes instruction-tuned models capable of following repair instructions---a prerequisite not previously documented.

\subsubsection{2.3 SMT-Guided Repair}

SMT-guided approaches enforce formal specification satisfaction. ContractEval \citep{lim2025contracteval} demonstrated that SMT solver integration achieves 75--82\% pass@1 on contract-satisfying assertions, compared to 0\% with standard prompting. Classical SMT-based repair tools such as Nopol and Angelix target human-written buggy code using Z3 constraint solving.

\subsubsection{2.4 Gap in Existing Work}

Each strategy stream has demonstrated effectiveness in isolation, but on different benchmarks with different metrics. No existing work measures whether these strategies target independent error classes, which is necessary to predict combined pipeline effectiveness.
